\section{Limitations and open problems}\label{sec:limits}

\paragraph{The $(\log q)^2$ barrier.}
Every prime-supported construction analysed through the explicit formula has guaranteed bias
$\asymp\log q/\sqrt X$. The reason is that about $(\log q)/2\pi$ zeros per unit height can conspire;
spreading the weight over scales changes this only by constants. An $\varepsilon$-biased prime set of size
$(\log q)^{2-\delta}$ for all $q$ would improve the GRH bound for the least prime outside a subgroup, which is
open \cite{lls2015,cmqr2024}. Montgomery's $\Omega(\log q\log\log q)$ \cite{montgomery1971} limits what is
possible even in principle. Empirically, typical moduli behave like random sets (Section~\ref{sec:experiments}),
with support $\asymp\log q\log X/\varepsilon^2$.

\paragraph{What is formal and what is not.}
Theorems~\ref{thm:A}, \ref{thm:sharp}, \ref{thm:closed} and Corollaries~\ref{cor:A}, \ref{cor:cert} are formal.
Proposition~\ref{prop:opt} is formal only in its abstract core (Lemma~\ref{lem:ranktwo}). Corollary~\ref{cor:asymp}
and Section~\ref{sec:complexity} are informal. Only one numerical instance of Theorem~\ref{thm:closed} is
machine-checked. The closed form loses a factor of about $4$ in support relative to the quadrature constants
(Table~\ref{tab:params}), mostly through the bound on $\mathrm{SP}(a)$ and the decay constants in $I_a(\beta)$.

\paragraph{Trust.}
Our unconditional theorems inherit the trust base of the Lean GRH proof. That proof has passed extensive
machine checks, but at the time of writing no mathematician has read it through, and its companion paper
does not contain a proof of the main inequality \cite{rubinsteinpage}.

\paragraph{Open problems.}
\begin{enumerate}[leftmargin=2em]
\item Formalize Proposition~\ref{prop:opt} and Corollary~\ref{cor:asymp}, and certify further instances,
  e.g.\ $q\le2^{1024}$ or $\varepsilon=10^{-2}$.
\item Beat the leading constant $c^*$, as a constant in front of $\log q/\sqrt X$, by leaving the positivity
  method. Can sign information on $\hat k$ beyond $\hat k\ge0$ help, as in the one-sided problem of \cite{cmqr2024}?
\item Extend to other groups whose characters are Dirichlet characters, such as quotients of $\Gq$ and abelian
  Galois groups over $\mathbb Q$. For nonabelian groups, e.g.\ $\mathrm{Aff}(\Z/q\Z)$, combine with zig-zag or
  semidirect-product arguments.
\item The class-group and isogeny-graph expanders of \cite{jmv2009} need GRH for Hecke $L$-functions of
  imaginary quadratic fields, which remains open.
\end{enumerate}

\paragraph{Acknowledgements and attribution.}
GRH for Dirichlet $L$-functions in Lean is due to Alpöge, with Claude \cite{alpoge2026,grhlean}. The explicit
formula library is \cite{zeta23}; the positivity and gamma-side lemmas and the test function are from
\cite{grhmiller}; the Perron formula is from \cite{schoenfeldlean}; the Borel--Carath\'eodory log-derivative
lemma is from \cite{rhprimeerror}. All are Apache-2.0. The Lean code, experiments and this text were developed
with the assistance of Claude (Anthropic).
